% Tutorial, Chapter 2 -- Building a Worksheet: Lessons 3 to 6 (AN-17).
% The running business: the Initech Systems Corporation supply room
% budget.  The keystrokes of each table are exactly those of the figure's
% trace (screens/tl3..tl6.k123); the figures of Lessons 4 to 6 start from
% the worksheet of the lesson before, saved by the shipped application from
% that lesson's keystrokes (samples/1-tu3.k123, 1-tu4.k123).
% Sources (Law 1): os/i123/app/face_a.c (entries: a value starts with a
% digit or + - . ( @ # $, an arrow key stores and moves; the label
% prefixes; /Copy with a single-cell TO; /File Save), os/i123/core
% (@SUM, recalculation), os/i123/fs/wk1.c (.WK1).
% GATES: test-flair-i123 test-flair-i123-files test-i123-formula test-i123-copy-move-model test-pubs-screens test-pubs-i123
\Gates{test-flair-i123, test-flair-i123-files, test-i123-formula,
  test-i123-copy-move-model, test-pubs-screens, test-pubs-i123}
\chapter{Building a Worksheet}

In this chapter you build a budget for the supply room of Initech Systems
Corporation: the cost of four supplies in each quarter of the year, with
totals.  You enter labels and numbers, write formulas, copy them, and save
the finished worksheet.  Start with an empty worksheet: if you have just
started Initech 123 you have one.

\section{Lesson 3: Entering Labels}
\ix{labels}
In this lesson you learn to:
\begin{itemize}
\item enter labels, the text of a worksheet
\item align a label at the right of its cell with a label prefix
\item enter several labels in a row or a column quickly
\end{itemize}

\subsection{The Title}
A \term{label} is an entry that begins with a letter, or with a label
prefix.  The cell pointer is in A1.

\begin{Keystrokes}
\Do{Type}{\type{INITECH SYSTEMS CORPORATION}}{}
\Do{Press}{\kname{DOWN}}{to enter the label and move to A2}
\Do{Type}{\type{Supply Room Budget}}{}
\Do{Press}{\kname{DOWN} twice}{to enter the label and move to A4}
\end{Keystrokes}
An arrow key enters what you typed and moves the cell pointer in one step.
The title is longer than column A is wide, so it runs on into the empty
cells to its right.

\subsection{The Column Headings}
Headings over numbers look best at the right of their cells, where the
numbers are.  A quotation mark typed before a label is a \term{label
prefix} that aligns the label at the right.

\begin{Keystrokes}
\Do{Type}{\type{Item}}{}
\Do{Press}{\kname{RIGHT}}{to move to B4}
\Do{Type}{\type{"Q1}}{the first quarter, right-aligned}
\Do{Press}{\kname{RIGHT}}{}
\Do{Type}{\type{"Q2}}{and press \kname{RIGHT}; then \type{"Q3}, \kname{RIGHT}, \type{"Q4}, \kname{RIGHT}}
\Do{Type}{\type{"Total}}{}
\Do{Press}{\kname{ENTER}}{to enter the last heading}
\end{Keystrokes}

\subsection{The Row Labels}
\begin{Keystrokes}
\Do{Press}{\fkey{Goto}{F5}}{and type \type{a5}, then press \kname{ENTER}}
\Do{Type}{\type{Paper}}{and press \kname{DOWN}}
\Do{Type}{\type{Toner}}{and press \kname{DOWN}}
\Do{Type}{\type{Staples}}{and press \kname{DOWN}}
\Do{Type}{\type{Folders}}{and press \kname{DOWN}}
\Do{Type}{\type{Total}}{and press \kname{ENTER}}
\end{Keystrokes}
Your worksheet now looks like Figure~\ref{fig:tl3}.
\FlairShot[fig:tl3]{tl3}{The labels of the supply room budget}

\section{Lesson 4: Entering Numbers}
\ix{values}
In this lesson you enter the cost of each supply in each quarter.  A
\term{value} is an entry that begins with a digit, or with one of the
characters \texttt{+ - . ( @ \# \$}.  Numbers align at the right of their
cells.

\begin{Keystrokes}
\Do{Press}{\fkey{Goto}{F5}}{and type \type{b5}, then press \kname{ENTER}}
\Do{Type}{\type{120}}{and press \kname{RIGHT}}
\Do{Type}{\type{135}}{and press \kname{RIGHT}}
\Do{Type}{\type{110}}{and press \kname{RIGHT}}
\Do{Type}{\type{150}}{and press \kname{ENTER}}
\end{Keystrokes}
Enter the other three rows in the same way, starting each at column B:

\begin{LotTable}{p{20mm}p{18mm}p{18mm}p{18mm}p{18mm}}
\LotHead{\textbf{Row} & \textbf{B} & \textbf{C} & \textbf{D} & \textbf{E}}
6 Toner & 210 & 190 & 240 & 205\\
7 Staples & 15 & 12 & 18 & 14\\
8 Folders & 40 & 35 & 45 & 30\LotEnd
\end{LotTable}
\FlairShot[fig:tl4]{tl4}{The numbers entered}

If you make a mistake in a number, move the cell pointer to it and type
the number again: the new entry replaces the old one.

\section{Lesson 5: Entering Formulas and Copying Them}
\ix{formulas}\ix{"@SUM}
In this lesson you learn to:
\begin{itemize}
\item enter a formula with the @SUM @function
\item copy a formula to other cells
\end{itemize}

\subsection{The Total of a Row}
A \term{formula} calculates a value.  The @function \fn{@SUM} adds the
numbers in a range: \type{@sum(b5..e5)} adds B5, C5, D5 and E5.

\begin{Keystrokes}
\Do{Press}{\fkey{Goto}{F5}}{and type \type{f5}, then press \kname{ENTER}}
\Do{Type}{\type{@sum(b5..e5)}}{the total of the paper row}
\Do{Press}{\kname{ENTER}}{}
\end{Keystrokes}
F5 shows the total, 515; the control panel shows the formula as Initech
123 stores it, \texttt{@SUM(B5..E5)} (Figure~\ref{fig:tl5a}).
\FlairShot[fig:tl5a]{tl5a}{The first total}

\subsection{Copying the Formula}
\ix{/Copy}
The same formula, one row lower each time, totals the other rows.  You do
not need to type it again: copy it.  /Copy copies the formula so that it
refers to the cells to its left in its new place.

\begin{Keystrokes}
\Do{Press}{\kname{SLASH}}{and type \type{c} to select Copy}
\Do{Press}{\kname{ENTER}}{to copy the current cell, F5}
\Do{Type}{\type{f6}}{the cell to copy to}
\Do{Press}{\kname{ENTER}}{}
\end{Keystrokes}
Now copy F5 and F6 together into F7 and F8.  Initech 123 places a copy at
the upper left cell you name, and the rest of the range follows it.

\begin{Keystrokes}
\Do{Press}{\kname{SLASH}}{and type \type{c}}
\Do{Type}{\type{f5..f6}}{the range to copy, and press \kname{ENTER}}
\Do{Type}{\type{f7}}{and press \kname{ENTER}}
\end{Keystrokes}

\subsection{The Totals of the Columns}
\begin{Keystrokes}
\Do{Press}{\fkey{Goto}{F5}}{and type \type{b9}, then press \kname{ENTER}}
\Do{Type}{\type{@sum(b5..b8)}}{and press \kname{ENTER}}
\Do{Press}{\kname{SLASH}}{type \type{c}, press \kname{ENTER}, type \type{c9}
  and press \kname{ENTER}}
\Do{Press}{\kname{SLASH}}{type \type{c}, type \type{b9..c9}, press
  \kname{ENTER}, type \type{d9} and press \kname{ENTER}}
\Do{Press}{\kname{SLASH}}{type \type{c}, type \type{b9..c9}, press
  \kname{ENTER}, type \type{e9} and press \kname{ENTER}}
\end{Keystrokes}
The last copy places the formulas of B9 and C9 in E9 and F9: F9 adds the
four totals above it (Figure~\ref{fig:tl5b}).
\FlairShot[fig:tl5b]{tl5b}{The finished budget}

\section{Lesson 6: Saving the Worksheet}
\ix{/File Save}\ix{saving}
The worksheet exists only in memory until you save it.  In this lesson you
save it in a file named SUPPLY.

\begin{Keystrokes}
\Do{Press}{\kname{SLASH}}{and type \type{f} to select File}
\Do{Type}{\type{s}}{to select Save}
\Do{Type}{\type{SUPPLY}}{the name of the file}
\Do{Press}{\kname{ENTER}}{to save the worksheet}
\end{Keystrokes}
Initech 123 adds the extension \file{.WK1} and saves the worksheet on drive
\file{A:}.  The status line at the bottom of the window shows
\texttt{Saved} and the file name (Figure~\ref{fig:tl6}).
\FlairShot[fig:tl6]{tl6}{The worksheet saved as SUPPLY}

This completes the chapter.  You can end Initech 123 now, with
\slashcmd{/Quit Yes}, or go on to the next chapter, which begins by
retrieving the file you saved.
